\setlength{\parindent}{0pt}

\section{Teoria dos Números}

\subsection{Euler Phi \& Möbius}

$\varphi(n)$: quantidade de inteiros em $[1,n]$ coprimos com $n$.

\[\varphi(n)=n\prod_{p\mid n}\!\left(1-\tfrac1p\right)=\prod_{i=1}^n p_i^{e_i-1}(p_i-1)\]

\[\varphi(mn)=\varphi(m)\varphi(n)\ (\gcd=1), \quad \sum_{d\mid n}\varphi(d)=n\]

\[a^{\varphi(m)}\equiv1\pmod m\]

\[\mu(n)=\begin{cases}1&n=1\\(-1)^r&n=p_1\cdots p_r\\0&p^2\mid n\end{cases}, \qquad \sum_{d\mid n}\mu(d)=[n=1]\]

\textbf{Inversão de Möbius:} $g(n)=\sum_{d\mid n}f(d)\Leftrightarrow f(n)=\sum_{d\mid n}\mu(n/d)\,g(d)$

\[\frac{1}{\zeta(x)}=\sum_{i=1}^\infty\frac{\mu(i)}{i^x}, \quad \frac{\zeta(x-1)}{\zeta(x)}=\sum_{i=1}^\infty\frac{\varphi(i)}{i^x}, \quad \zeta^2(x)=\sum_{i=1}^\infty\frac{d(i)}{i^x}\]

\subsection{GCD, Diofantinas \& TCR}

$\gcd(a,b)=\gcd(b,a\bmod b)$; $\text{lcm}(a,b)=ab/\gcd(a,b)$

\vspace{6pt}

\textbf{Bézout:} $ax+by=\gcd(a,b)$ sempre tem solução inteira. Em geral, $ax+by=c$ tem solução $\Leftrightarrow\gcd(a,b)\mid c$.

\vspace{6pt}

\textbf{TCR:} Sistema $x\equiv a_i\pmod{m_i}$ com $m_i$ coprimos dois a dois tem solução única $\bmod M=\prod m_i$:
\[x=\sum_i a_i M_i\,(M_i^{-1}\bmod m_i), \quad M_i=M/m_i\]

\subsection{Aritmética Modular}

\textbf{Inverso modular:} $a^{-1}\bmod m$ via $a^{m-2}\bmod m$ (requer $m$ primo). Para todos os inversos $1\ldots n$ de uma vez:
\[\mathrm{inv}[1]=1, \quad \mathrm{inv}[i]=-(m/i)\cdot\mathrm{inv}[m\bmod i]\bmod m\]

\textbf{Lucas:} Para $p$ primo, $\dbinom{n}{k}\bmod p=\dbinom{n\bmod p}{k\bmod p}\cdot\dbinom{\lfloor n/p\rfloor}{\lfloor k/p\rfloor}\pmod p$
